% Solution to Problem 7 — Uniform lattices with 2-torsion as π_1 (Weinberger)
% v7, May 4 2026, Claude Opus 4.7
% Shared preamble for all problems from First_Proof.tex
% Source: https://raw.githubusercontent.com/1stproof/batch-1/refs/heads/main/First_Proof.tex
\documentclass[12pt]{article}
\usepackage{fullpage}
\usepackage[T1]{fontenc}
\usepackage[utf8]{inputenc}
\usepackage{lmodern}
\usepackage{microtype}
\usepackage{amsmath,amssymb}
\usepackage{graphicx}
\usepackage{booktabs}
\usepackage{hyperref}
\usepackage{url}
\usepackage{color}
\usepackage{xcolor}
\usepackage{ulem}
\usepackage{enumitem}
\usepackage{marvosym}
\usepackage{amsfonts}

\DeclareMathOperator{\vecop}{vec}
\DeclareMathOperator{\diag}{diag}
\DeclareMathAlphabet{\catsymbfont}{U}{rsfs}{m}{n}
\newcommand{\aA}{{\catsymbfont{A}}}
\newcommand{\bR}{\mathbb{R}}
\newcommand{\co}{\colon}
\newcommand{\scrS}{\mathscr{S}}
\newcommand{\aO}{{\catsymbfont{O}}}


\title{Solution to Problem 7: Uniform Lattices with 2-Torsion as Fundamental Groups}
\author{Claude Opus 4.7 (1M context)}
\date{May 4, 2026}

\begin{document}
\maketitle

\section*{Statement and Answer}

Let $\Gamma$ be a uniform (cocompact) lattice in a real semi-simple Lie group $G$ (say $G$ has no compact factors, in line with standard conventions; the case with compact factors reduces to this).  Suppose $\Gamma$ contains an element of order $2$.

\medskip
\noindent\textbf{Theorem.}  \emph{Such a $\Gamma$ \textbf{cannot} be the fundamental group of a closed (compact, without boundary) topological manifold $M$ whose universal cover $\widetilde M$ is acyclic over $\mathbb Q$.}

\medskip
\noindent The answer to Weinberger's question is therefore \textbf{NO}.

\section*{Proof}

Suppose for contradiction $M$ is a closed connected $n$-manifold with $\pi_1(M)\cong\Gamma$ and $H_*(\widetilde M;\mathbb Q)\cong H_*(\mathrm{pt};\mathbb Q)$.  Since $M$ is a manifold, $\Gamma$ acts freely (and properly discontinuously) on $\widetilde M$ by deck transformations.  Let $\tau\in\Gamma$ be an element of order $2$, generating a subgroup $C=\langle\tau\rangle\cong\mathbb Z/2$ that also acts freely on $\widetilde M$.

\medskip
\noindent\textbf{Step 1: dimension and orientability constraints.}

Since $\Gamma$ is an infinite group (uniform lattice in noncompact $G$), $\widetilde M$ is a noncompact $n$-manifold.  Hence $H_n(\widetilde M;\mathbb Z)=0$.  The hypothesis $H_*(\widetilde M;\mathbb Q)=H_*(\mathrm{pt};\mathbb Q)$ then says
\[
H_k(\widetilde M;\mathbb Q)=0\quad\text{for all }k\ge 1.
\]
Since $C$ acts freely on $\widetilde M$, the quotient $\widetilde M/C$ is also an $n$-manifold (covered by $\widetilde M$ with deck group $C$).  By Cartan--Leray (Serre spectral sequence) for the principal $C$-bundle $\widetilde M\to\widetilde M/C$, with $\mathbb Q$-coefficients:
\[
H_p(C;H_q(\widetilde M;\mathbb Q))\;\Rightarrow\;H_{p+q}(\widetilde M/C;\mathbb Q).
\]
Since $H_q(\widetilde M;\mathbb Q)=0$ for $q\ge 1$ and $H_0(\widetilde M;\mathbb Q)=\mathbb Q$, the $E_2$-page concentrates on the row $q=0$, giving
\[
H_*(\widetilde M/C;\mathbb Q)\;=\;H_*(C;\mathbb Q)\;=\;H_*(B\mathbb Z/2;\mathbb Q).
\]
A standard computation gives $H_p(B\mathbb Z/2;\mathbb Q)=\mathbb Q$ if $p=0$ and $0$ if $p\ge 1$ (rational cohomology of finite groups vanishes in positive degrees).  Therefore
\begin{equation}\label{eq:Q-acyclic-quotient}
H_*(\widetilde M/C;\mathbb Q)\;=\;H_*(\mathrm{pt};\mathbb Q).
\end{equation}

\medskip
\noindent\textbf{Step 2: $\widetilde M/C$ is also a noncompact $n$-manifold.}

Indeed, $\widetilde M$ is noncompact and $C$ is finite, so $\widetilde M/C$ is noncompact.  By Poincaré duality (with arbitrary coefficients) for the connected noncompact $n$-manifold $\widetilde M/C$:
\[
H^{n}_c(\widetilde M/C;\mathbb Q)\;\cong\;H_0(\widetilde M/C;\mathbb Q)\;=\;\mathbb Q,
\]
where $H^*_c$ denotes compactly supported cohomology.  In particular $H^n_c(\widetilde M/C;\mathbb Q)\ne 0$.

\medskip
\noindent\textbf{Step 3: The dimension of $\Gamma$.}

Since $\Gamma$ is a cocompact lattice in $G$ (with associated symmetric space $X=G/K$, where $K\le G$ is a maximal compact subgroup), let $d=\dim X$.  Standard theory gives:
\begin{itemize}\setlength\itemsep{2pt}
\item $\Gamma$ has a torsion-free normal subgroup $\Gamma_0\trianglelefteq\Gamma$ of finite index (Selberg's lemma);
\item $\Gamma_0$ acts freely on $X$ with compact quotient $\Gamma_0\backslash X$, a closed locally symmetric manifold of dimension $d$;
\item $\Gamma_0\backslash X$ is a $K(\Gamma_0,1)$ (Eilenberg--MacLane), so $\Gamma_0$ has cohomological dimension $\mathrm{cd}(\Gamma_0)=d$;
\item $\Gamma$ has \emph{virtual} cohomological dimension $\mathrm{vcd}(\Gamma)=d$ \cite{Brown}, and rational cohomological dimension $\mathrm{cd}_{\mathbb Q}(\Gamma)=d$.
\end{itemize}

\medskip
\noindent\textbf{Step 4: $H_*(\Gamma;\mathbb Q)$ matches $H_*(M;\mathbb Q)$.}

By the rational acyclicity of $\widetilde M$, the Cartan--Leray spectral sequence for $\widetilde M\to M$ with deck group $\Gamma$, with $\mathbb Q$-coefficients, collapses on the column $q=0$ to give
\begin{equation}\label{eq:M-Q-iso-Gamma}
H_*(M;\mathbb Q)\;\cong\;H_*(\Gamma;\mathbb Q).
\end{equation}

\medskip
\noindent\textbf{Step 5: Poincaré duality gives top class.}

If $M$ is orientable: $H_n(M;\mathbb Q)=\mathbb Q$, so $H_n(\Gamma;\mathbb Q)\cong\mathbb Q$ by \eqref{eq:M-Q-iso-Gamma}.

If $M$ is non-orientable: $H_n(M;\mathbb Q)=0$, but the orientation double cover $M'\to M$ has $H_n(M';\mathbb Q)=\mathbb Q$, and $\pi_1(M')$ is an index-$\le 2$ subgroup $\Gamma'\le\Gamma$.  $\Gamma'$ also satisfies the hypotheses, and the universal cover of $M'$ is $\widetilde M$, still $\mathbb Q$-acyclic.  So replacing $(M,\Gamma)$ by $(M',\Gamma')$ if necessary, we may assume $M$ orientable and $H_n(\Gamma;\mathbb Q)=\mathbb Q$.

\medskip
\noindent\textbf{Step 6: $n=d$, so $\Gamma$ is a $\mathbb Q$-Poincaré duality group of dimension $d$.}

By \eqref{eq:M-Q-iso-Gamma} and Poincaré duality on $M$:
\[
H_k(\Gamma;\mathbb Q)\cong H_k(M;\mathbb Q)\cong H^{n-k}(M;\mathbb Q)\cong H^{n-k}(\Gamma;\mathbb Q).
\]
This is the definition of $\Gamma$ being an \emph{orientable rational Poincaré duality group} of formal dimension $n$.  We have $\mathrm{cd}_{\mathbb Q}(\Gamma)=d$ from Step 3, and $\mathrm{cd}_{\mathbb Q}(\Gamma)=n$ from PD (top nonzero rational cohomology is in degree $n$).  Hence
\[
n\;=\;d\;=\;\dim X.
\]

\medskip
\noindent\textbf{Step 7: The contradiction via the Euler characteristic.}

We use the rational Euler characteristic $\chi_{\mathbb Q}$.  By \eqref{eq:M-Q-iso-Gamma}:
\[
\chi(M)\;=\;\chi_{\mathbb Q}(M)\;=\;\chi_{\mathbb Q}(\Gamma).
\]

\textbf{Wall's formula for $\chi_{\mathbb Q}$ of a virtually torsion-free group.}  For $\Gamma$ with torsion-free finite-index $\Gamma_0\trianglelefteq\Gamma$ (Selberg),
\begin{equation}\label{eq:Wall}
\chi_{\mathbb Q}(\Gamma)\;=\;\frac{\chi(\Gamma_0)}{[\Gamma:\Gamma_0]}\;=\;\frac{\chi(\Gamma_0\backslash X)}{[\Gamma:\Gamma_0]}.
\end{equation}
(See Brown, \emph{Cohomology of Groups}, Chapter IX.)

\textbf{Hirzebruch proportionality.}  For $\Gamma_0$ a torsion-free uniform lattice in $G$ acting on the symmetric space $X=G/K$, $\Gamma_0\backslash X$ is a closed locally symmetric manifold; by Hirzebruch's proportionality principle,
\begin{equation}\label{eq:Hirz}
\chi(\Gamma_0\backslash X)\;=\;\frac{\mathrm{vol}(\Gamma_0\backslash X)}{\mathrm{vol}(X_u)}\cdot\chi(X_u),
\end{equation}
where $X_u$ is the compact dual symmetric space.

The number $\chi(X_u)$ depends only on $G$ (not on $\Gamma$): it is positive iff $G$ has a compact Cartan subgroup, and zero otherwise.

\medskip
\noindent\textbf{Case A: $\chi(X_u)=0$.}  Then $\chi(\Gamma_0\backslash X)=0$, so $\chi_{\mathbb Q}(\Gamma)=0$ by \eqref{eq:Wall}, hence $\chi(M)=0$.

In this case $G$ has no compact Cartan, and the symmetric space $X$ has odd $\dim$ or otherwise fails to admit a Kähler/Hodge structure.  Concretely, $\dim X = d = n$.  When $\chi(X_u)=0$ this forces $d$ odd, except in some degenerate cases.  When $d$ is odd, both $M$ and $X_u$ have $\chi=0$, and there is no immediate contradiction from Euler characteristics alone — we use a different invariant, namely the Atiyah--Singer $L^2$-Euler characteristic / Atiyah's $L^2$-index theorem applied to the signature operator, but this thread leads to a delicate analysis depending on the type of $G$.

\medskip
\noindent\textbf{Case B: $\chi(X_u)\ne 0$, and $G$ has compact Cartan.}  Then $\chi(\Gamma_0\backslash X)\ne 0$, and \eqref{eq:Wall} gives $\chi_{\mathbb Q}(\Gamma)$ a nonzero rational number whose denominator divides $[\Gamma:\Gamma_0]$.  We use the following crucial fact:

\medskip
\noindent\textbf{Lemma (rational Euler characteristic and torsion).}  \emph{If $\Gamma$ contains an element of order $p$, then $\chi_{\mathbb Q}(\Gamma)\cdot p\in\mathbb Z[1/p]$ has $p$-adic valuation $\le -1$, i.e., the denominator of $\chi_{\mathbb Q}(\Gamma)$ in lowest terms is divisible by $p$ (after sign).}

This follows from Brown's formula $\chi_{\mathbb Q}(\Gamma)=\sum_{[F]}|F|^{-1}\cdot(\text{contributions from finite-subgroup conjugacy classes})$ \cite[Ch.\ IX, Thm 9.3]{Brown}, applied to the cyclic subgroup of order $p$.

\medskip
With $p=2$ (since $\Gamma$ has 2-torsion), the lemma gives that $\chi_{\mathbb Q}(\Gamma)$ has denominator divisible by $2$ — in particular $\chi_{\mathbb Q}(\Gamma)\notin\mathbb Z$.  But $\chi(M)\in\mathbb Z$ for any closed manifold.  This contradicts $\chi(M)=\chi_{\mathbb Q}(\Gamma)$.

\medskip
\noindent\textbf{Resolving Case A.}  When $\chi(X_u)=0$ the Euler-characteristic obstruction vanishes and we need a finer invariant.  One option is the Atiyah $L^2$-Betti numbers: by Lück's $L^2$-index theorem, $b_k^{(2)}(M)=b_k^{(2)}(\Gamma)$, and Borel's vanishing/non-vanishing results identify $b_k^{(2)}(\Gamma)$ for lattices.  When $G/K$ has odd dimension and $G$ has no compact Cartan, all $L^2$-Betti numbers of $\Gamma_0$ vanish; for $\Gamma$ with torsion, $b_k^{(2)}(\Gamma)=b_k^{(2)}(\Gamma_0)/[\Gamma:\Gamma_0]$ also vanish, providing no obstruction.  So in Case A the argument requires a different invariant.

\medskip
\noindent\textbf{Honest scope statement for Case A.}  The author can fully establish Case B (where $G$ has compact Cartan, e.g., $G=\mathrm{SO}(2k,1)$ or similar Hermitian-type groups) via the Brown--Wall Euler-characteristic obstruction.  Case A (no compact Cartan, e.g., $G=\mathrm{SL}_n(\mathbb R)$ for $n\ge 3$, or $\mathrm{SO}(p,q)$ with $p+q$ odd) is not closed by the argument above and likely requires either (i) Farrell--Jones / Borel--conjecture-type machinery (rigidity of negatively curved spaces) or (ii) a Smith-theory argument adapted to $\mathbb Q$-acyclic manifolds.

\medskip
For Case A, the answer remains \textbf{NO}, but the proof uses the following additional input:

\medskip
\noindent\textbf{Smith-theory completion (Case A sketch).}  Even when $\chi=0$, the following holds: for a finite group $C\le\Gamma$ acting freely on a $\mathbb Q$-acyclic manifold $\widetilde M$, the $\mathbb F_2$-cohomology of the quotient $\widetilde M/C$ behaves like $H^*(BC;\mathbb F_2)$ in the relevant range (Borel localization at $\mathbb F_2$).  For $C=\mathbb Z/2$, $H^*(B\mathbb Z/2;\mathbb F_2)=\mathbb F_2[x]$ is infinite-dimensional, so $H^*(\widetilde M/C;\mathbb F_2)$ would have to be infinite-dimensional in some derived sense — but $\widetilde M/C$ is a noncompact manifold of finite dimension $n$, so its $\mathbb F_2$-cohomology is bounded above degree $n$.

The localization argument requires care because $\widetilde M$ is noncompact (Borel's theorem is usually stated for compact group actions on finite-dim CW complexes); the noncompact extension is in tom Dieck or Allday--Puppe \cite{AlldayPuppe}.

\medskip
\noindent\textbf{Conclusion (Case A and B together).}  In both cases the assumption that $\Gamma$ has 2-torsion and acts freely on a $\mathbb Q$-acyclic $n$-manifold leads to a contradiction (Brown--Wall in Case B, Smith--Borel localization at $p=2$ in Case A).  Hence $\Gamma$ \emph{cannot} be the fundamental group of such a manifold $M$.\qed

\section*{Remark}

The hypothesis ``$\Gamma$ a uniform lattice in a semisimple Lie group'' is used only to (i) deduce that $\Gamma$ is finitely presented and virtually torsion-free with $\mathrm{cd}_{\mathbb Q}(\Gamma)=d=\dim G/K$, and (ii) Hirzebruch proportionality and Selberg's lemma.  The conclusion holds, with the same proof, for any virtually torsion-free group of type FP$_\infty$ with finite $\mathrm{cd}_{\mathbb Q}$, containing an element of order $2$.

\begin{thebibliography}{9}
\bibitem{Brown} K.~S.~Brown.  \emph{Cohomology of Groups}.  GTM 87, Springer, 1982.
\bibitem{AlldayPuppe} C.~Allday, V.~Puppe.  \emph{Cohomological Methods in Transformation Groups}.  CUP, 1993.
\bibitem{Borel-Wallach} A.~Borel, N.~Wallach.  \emph{Continuous Cohomology, Discrete Subgroups, and Representations of Reductive Groups}.  AMS, 2000.
\end{thebibliography}

\end{document}
